
Each system component is validated through targeted experiments testing specific mechanistic claims. Experiments use proof-of-concept scale infrastructure with mock or simulated data to isolate mechanism feasibility from deployment complexity.

\subsubsection{4.1 Experimental Questions}

\begin{enumerate}
\item \textbf{Q1 (h-e1):} Can load-time instrumentation track adoption with <10\% overhead and $\geq 95$\% capture rate?
\item \textbf{Q2 (h-m1):} Do health metrics achieve $\geq 60$\% precision and $\geq 80$\% recall in identifying deprecation candidates?
\item \textbf{Q3 (h-m2):} Does context-aware inference achieve $\geq 70$\% accuracy on task-specific successors?
\item \textbf{Q4 (h-m3):} Can dependency graphs achieve $\geq 95$\% impact completeness for migration planning?
\item \textbf{Q5 (h-m4):} Does the full telemetry stack maintain $\geq 95$\% capture with <10\% overhead?
\end{enumerate}

\subsubsection{4.2 Datasets and Baselines}

\textbf{HuggingFace Datasets Hub Infrastructure:}
\begin{itemize}
\item \textbf{h-e1:} 100 benchmark dataset loads (CIFAR-10, IMDB, WikiText-103) to measure overhead across 10 runs per dataset.
\item \textbf{h-m1:} 1000-dataset simulation with synthetic health metrics (velocity, emergence, issue\_ratio) over a 6-month observation window.
\item \textbf{h-m2:} 500 synthetic validation cases (70/30 train/test split) with ground-truth task labels across four contexts (classification, pretraining, robustness, unknown).
\item \textbf{h-m3:} Mock 4-node dependency graph (deprecated dataset $\rightarrow$ successor dataset, dependent model, dependent pipeline) with schema compatibility checks.
\item \textbf{h-m4:} Integrated stack with 100 simulated user sessions (4 loads per user, 40\% deprecated dataset encounter rate, 50\% adoption rate).
\end{itemize}

\textbf{Baselines:} No formal baselines exist because no prior dataset deprecation systems exist. Performance targets are derived from software package manager requirements (NPM overhead benchmarks, PyPI deprecation warning thresholds).

\textbf{Mock Data Rationale:} Real HuggingFace API integration was blocked by configuration errors and rate limits during proof-of-concept validation. Mock data isolates mechanism validation from external dependencies. Generalization to real API data is future work.

\subsubsection{4.3 Evaluation Metrics}

\begin{itemize}
\item \textbf{Overhead:} \texttt{(instrumented\_time - baseline\_time) / baseline\_time $\times$ 100\%}
\item \textbf{Capture Rate:} \texttt{captured\_events / total\_events}
\item \textbf{Precision:} \texttt{true\_positives / (true\_positives + false\_positives)}
\item \textbf{Recall:} \texttt{true\_positives / (true\_positives + false\_negatives)}
\item \textbf{Context Accuracy:} \texttt{correct\_inferences / total\_inferences}
\item \textbf{Override Rate:} \texttt{user\_overrides / total\_recommendations}
\end{itemize}

Wilson confidence intervals (95\% CI) were computed for capture rate to validate statistical reliability.

\subsubsection{4.4 Experimental Protocol}

Each sub-hypothesis follows a gate protocol:
\begin{enumerate}
\item \textbf{Setup:} Load mock or simulated data matching the experimental question.
\item \textbf{Execution:} Run mechanism under test conditions.
\item \textbf{Measurement:} Compute target metrics (overhead, precision, accuracy).
\item \textbf{Gate Check:} Pass if all metrics meet thresholds, otherwise fail.
\end{enumerate}

MUST\_WORK gates (h-e1, h-m1, h-m2) are critical for system viability. SHOULD\_WORK gates (h-m3, h-m4) validate supporting infrastructure but do not block progression.

\subsubsection{4.5 Reproducibility}

Code, mock data generators, and evaluation scripts are available in the research directory under \texttt{h-e1/}, \texttt{h-m1/}, \texttt{h-m2/}, \texttt{h-m3/}, and \texttt{h-m4/}. Simulated HuggingFace metadata uses patterns from literature on software package deprecation velocity distributions and GitHub issue accumulation rates. Proof-of-concept scale was chosen to validate mechanism feasibility; production validation requires deployment testing.
